\chapter{Conclusion}
\label{ch:conclusion}

This thesis set out to answer whether a purely statistical, interpretable model
could predict impending Sepsis-3 onset ahead of clinical recognition on
MIMIC-IV, matching or exceeding established rule-based scores and surviving
external validation. The motivation was a specific imbalance in the literature:
statistical work on MIMIC-IV overwhelmingly predicts mortality after diagnosis,
rarely tests the proportional-hazards assumption, and seldom validates
onset-prediction models externally.

The work delivered an interpretable pipeline whose central deliverable is a
dynamic landmark supermodel (C1): a single elastic-net logistic model fitted on
207{,}010 stacked landmark rows across 48{,}829 stays, producing a pooled
out-of-sample AUROC of \textbf{0.797} across six prediction hours. The
supermodel matches or exceeds bespoke per-landmark models while remaining one
parsimonious, deployable object that emits continuously updated hourly risk.
The pipeline formally tests and acts on the proportional-hazards assumption
(C2), reframes the benchmark by showing that bedside severity scores fail on
the incident-onset task while a multivariable nomogram reaches an AUROC of
0.775 (C3), characterises an informative-missingness ceiling in which whether
a laboratory test was ordered outweighs its value (C4), and, closing the
validation gap, transports the frozen model to eICU-CRD (C5). External
discrimination held at 0.686 against an internal eICU ceiling of 0.711, and
the expected miscalibration was corrected by a single-parameter recalibration.

The central lesson is that, for retrospectively labelled ICU data, the ceiling
on onset prediction is set by the labelling process rather than by model
complexity; an interpretable statistical model is therefore a scientifically
honest and clinically transparent choice. Avenues for further research follow
naturally. The same landmark substrate can be used to benchmark machine-learning
approaches against this interpretable baseline on identical splits; the supermodel
can be transported across all landmark hours rather than the hour-six landmark
alone; and prospective, multi-centre evaluation would test whether the
recalibrated model supports clinically useful early warning at an acceptable
alarm burden. By pairing rigorous internal methodology with a genuine external
test, this work offers a reproducible and transportable baseline for the next
generation of sepsis early-warning research.
